\documentclass{training}
\schoollogo{../../../resources/latex/crack.pdf}
\usepackage[utf8]{inputenc}
\usepackage[T1]{fontenc}
\usepackage{lmodern}
\setlang{french}
\usepackage{hyperref}
\usepackage{graphicx}
\usepackage{tabularx}
\usepackage{fvextra}
\fvset{breaklines=true,breakanywhere=true,breaksymbol={},breaksymbolleft={},breaksymbolright={},breakanywheresymbolpre={},breakanywheresymbolpost={}}
\makeatletter
\g@addto@macro\UrlBreaks{%
  \do\a\do\b\do\c\do\d\do\e\do\f\do\g\do\h\do\i\do\j\do\k\do\l\do\m%
  \do\n\do\o\do\p\do\q\do\r\do\s\do\t\do\u\do\v\do\w\do\x\do\y\do\z%
  \do\A\do\B\do\C\do\D\do\E\do\F\do\G\do\H\do\I\do\J\do\K\do\L\do\M%
  \do\N\do\O\do\P\do\Q\do\R\do\S\do\T\do\U\do\V\do\W\do\X\do\Y\do\Z%
  \do\0\do\1\do\2\do\3\do\4\do\5\do\6\do\7\do\8\do\9%
  \do\<\do\>\do\(\do\)\do\'\do\/}
\schoollogo{../../../resources/latex/crack.pdf}
\makeatother


\begin{document}

\title{Le modèle de mutation --- M-EPI-001}
\subtitle{Familles, chemins et hypothèses}
\member{Ted ``Tedsig42'' Kouhouenou}{ted.sig42@tutamail.com}
\summary{Classer un composant dans la bonne famille architecturale, distinguer le chemin nominal des chemins d'abandon, et reconnaître les deux modes de mutation décrits par le papier : incomplete-fix variant et architectural scatter.}
\version{1.0}
\prerequisite{M-EPI-000, bases de C et de Linux}
\maketitle

\newpage
\tableofcontents
\newpage

\chapter{Instructions}
\begin{itemize}
  \item Une journée, en binôme de préférence.
  \item Les exercices sont des analyses : pas de code d'exploitation.
  \item Pour chaque classification, écrivez la question qui a tranché.
  \item Citez les sources et assumez les incertitudes.
\end{itemize}

\subsection*{Format du rendu --- en-tête obligatoire}

Tout fichier rendu dans ce module doit commencer par l'en-tête
ci-dessous, \textbf{adapté à la syntaxe de commentaire du langage}
utilisé.

\subsubsection*{Modèle générique}

\begin{minted}{c}
/**
 * Module: <code module>
 * Exercise: <numéro exo>
 * File: <nom du fichier>
 * Author: <votre login JeuneCrack>
 *
 * Description: <une phrase, pas de roman>
 *
 * Return: <voir ci-dessous>
 */
\end{minted}

\subsubsection*{Login}

Votre identifiant \textbf{JeuneCrack} (ex: \texttt{tedsig42}).
Pour un rendu en binôme : \texttt{login1/login2}.
Ajoutez un handle Matrix après le login si vous le souhaitez
(ex: \texttt{tedsig42 @tedsig42}).

\subsubsection*{Return --- selon le type de fichier}

\begin{tabularx}{\linewidth}{lX}
  \toprule
  \textbf{Type de fichier} & \textbf{Valeur de Return} \\
  \midrule
  C (\texttt{.c}) &
  \texttt{return (0);} (succès),
  \texttt{return (-1);} (erreur),
  \texttt{return (valeur);} (résultat calculé) \\

  Python (\texttt{.py}) &
  \texttt{0} (succès),
  \texttt{1} (erreur),
  \texttt{N/A} (analyse, script sans exit) \\

  Shell (\texttt{.sh}) &
  \texttt{exit 0} (succès),
  \texttt{exit 1} (erreur) \\

  Markdown (\texttt{.md}) &
  \texttt{N/A} --- les fichiers d'analyse n'ont pas de valeur de retour \\

  Exploit C/Python &
  \texttt{shell} (shell obtenu),
  \texttt{flag: <flag>} (flag capturé),
  \texttt{proof: <hash>} (preuve de concept) \\

  Binaire (\texttt{.bin}) &
  \texttt{N/A} --- pas d'en-tête pour les fichiers binaires \\
  \bottomrule
\end{tabularx}

\subsubsection*{Exemples par langage}

\textbf{C (\texttt{.c})} :
\begin{minted}{c}
/**
 * Module: F-EXP-001
 * Exercise: ex02
 * File: exploit.c
 * Author: tedsig42
 *
 * Description: Buffer overflow with ret2libc
 *
 * Return: 0
 */
\end{minted}

\textbf{Python (\texttt{.py})} :
\begin{minted}{python}
# Module: F-WEB-001
# Exercise: ex05
# File: ssrf-exploit.py
# Author: tedsig42
#
# Description: Blind SSRF bypass via redirect
#
# Return: 0
\end{minted}

\textbf{Shell (\texttt{.sh})} :
\begin{minted}{bash}
# Module: F-SUP-001
# Exercise: ex03
# File: dependency-confusion.sh
# Author: tedsig42
#
# Description: Automated dependency confusion test
#
# Return: exit 0
\end{minted}

\textbf{Markdown (\texttt{.md})} --- analyse, rapport :
\begin{Verbatim}[fontsize=\normalsize]
<!--
Module: M-SRC-001
Exercise: ex04
File: audit-report.md
Author: tedsig42

Description: Code audit of vulnerable-httpd

Return: N/A
-->
\end{Verbatim}

\chapter{Instructions Communes}

\section{Structure du D\'ep\^ot}

Votre d\'ep\^ot Git doit respecter la convention de nommage suivante pour
tous les modules soumis. L'outil d'\'evaluation (\texttt{pendragon})
se base sur cette structure pour localiser et tester votre travail.

\begin{console}
<login>/
+-- <MODULE-ID>-<slug>/         # ex. M-BLS-001-business-logic/
|   +-- <turnindir>/            # ex. ex00-pricing-manipulation/
|   |   +-- <file1>             # ex. poc-pricing.py
|   |   +-- <file2>             # ex. report.md
|   |   +-- ...
|   +-- <turnindir>/            # ex. ex01-race-conditions/
|   |   +-- ...
|   +-- ...
+-- ...
\end{console}

\begin{itemize}
  \item \textbf{R\'epertoire du module} : le nom complet du r\'epertoire tel
        qu'affich\'e dans votre curriculum (ex. \texttt{M-BLS-001-business-logic}).
  \item \textbf{R\'epertoire de rendu} : la valeur de \texttt{\textbackslash turnindir}
        fournie dans chaque exercice (ex. \texttt{ex00-pricing-manipulation}).
  \item \textbf{Fichiers de l'exercice} : exactement les fichiers list\'es dans
        \texttt{\textbackslash exfiles}. Ni fichiers suppl\'ementaires, ni fichiers manquants.
\end{itemize}

\section{Qu'est-ce qu'une Preuve de Concept Valide}

Une PoC doit \^etre :

\begin{enumerate}
  \item \textbf{Reproductible}. N'importe qui avec le m\^eme environnement peut
        l'ex\'ecuter et observer la vuln\'erabilit\'e. Les chemins cod\'es en dur,
        les adresses IP absolues et les jetons expir\'es ne sont pas acceptables.
  \item \textbf{Document\'ee}. Chaque \'etape est expliqu\'ee. L'\'evaluteur ne
        devrait pas avoir \`a deviner ce que fait le script.
  \item \textbf{Cibl\'ee}. L'exercice sp\'ecifie une cible autoris\'ee
        (lab, maquette, binaire local ou plateforme d\'esign\'ee). Ne testez
        jamais les PoC sur des syst\`emes de production qui ne vous appartiennent pas.
  \item \textbf{Minimale}. La PoC d\'emontre \emph{uniquement} la vuln\'erabilit\'e
        en question. Les charges utiles non li\'ees et le bruit sont p\'enalis\'es.
\end{enumerate}

\begin{warning}
  Une PoC qui modifie silencieusement l'\'etat du syst\`eme, supprime des
  fichiers ou persiste des changements en dehors du p\'erim\`etre de l'exercice
  sera rejet\'ee. Si votre exploit n\'ecessite un nettoyage, documentez les
  \'etapes de nettoyage dans votre rapport.
\end{warning}

\section{Processus de Soumission}

\begin{enumerate}
  \item R\'ealisez les livrables de l'exercice sur votre machine locale.
  \item Committez et poussez vers votre d\'ep\^ot Git assign\'e.
  \item Le d\'ep\^ot est gel\'e apr\`es la date limite.
  \item Les pairs clonent votre d\'ep\^ot et \'evaluent en utilisant :
        \begin{itemize}
          \item \texttt{pendragon} pour l'ex\'ecution automatis\'ee des tests
                (l\`a o\`u le module fournit des configurations mouli).
          \item La grille manuelle pour la qualit\'e du rapport, le style de
                code et la m\'ethodologie.
        \end{itemize}
\end{enumerate}

\hint{Ex\'ecutez toujours \texttt{pendragon run <MODULE-ID> <ex\_id>
<path/to/student/repo>} avant votre commit final pour v\'erifier
que votre soumission passe les v\'erifications automatis\'ees.}

\subsection*{Test avec pendragon}

Si le module inclut une configuration mouli (fichiers JSON dans
\texttt{resources/mouli/}), vous pouvez tester votre soumission localement :

\begin{console}
pendragon run M-BLS-001 ex00 ~/student-repos/tedsig42
\end{console}

\texttt{pendragon} r\'esout le chemin de l'exercice automatiquement en
cherchant \texttt{<mod-fullname>/<ex-dirname>/} dans le r\'epertoire
\'etudiant donn\'e.

\subsection*{V\'erification de style avec caliban}

Pour les modules qui n\'ecessitent du code C, \texttt{caliban} v\'erifie
la conformit\'e du style :

\begin{console}
caliban check ~/student-repos/tedsig42/M-BLS-001-business-logic
\end{console}

\texttt{caliban} applique la Norme 42 (taille des fonctions, nommage
des variables, constructions interdites) adapt\'ee aux contextes de
recherche de vuln\'erabilit\'es. Ex\'ecutez-le avant la soumission ---
un \'echec aux v\'erifications de style peut invalider l'exercice.

\section{R\`egle des Fichiers Suppl\'ementaires Interdits}

\begin{itemize}
  \item Votre d\'ep\^ot doit contenir \textbf{uniquement} les fichiers
        list\'es dans \texttt{\textbackslash exfiles} pour chaque exercice.
  \item Fichiers interdits : \texttt{.DS\_Store}, \texttt{\_\_pycache\_\_},
        \texttt{.venv/}, \texttt{node\_modules/}, \texttt{.git/},
        binaires compil\'es (sauf explicitement list\'es dans \texttt{exfiles}),
        fichiers d'\'echange d'\'editeur, fichiers journaux et tout fichier
        non demand\'e.
  \item Utilisez \texttt{.gitignore} pour exclure automatiquement les
        fichiers ind\'esirables.
\end{itemize}

\section{Format du Rapport}

Lorsqu'un rapport (\texttt{report.md}) est requis, il doit inclure :

\begin{enumerate}
  \item \textbf{Classe de vuln\'erabilit\'e} (identifiant CWE si applicable).
  \item \textbf{Étapes de reproduction} --- num\'erot\'ees, pr\^etes
        \`a \^etre copi\'ees-coll\'ees.
  \item \textbf{Charges utiles} --- chaînes exactes, requ\^etes ou
        entr\'ees utilis\'ees.
  \item \textbf{Impact} --- ce qu'un attaquant peut accomplir.
  \item \textbf{Mitigation} --- comment le d\'eveloppeur devrait corriger.
  \item \textbf{Sources} --- toutes les ressources consult\'ees (voir
        Section~\ref{sec:cite-sources}).
\end{enumerate}

\subsection*{En-t\^ete de Rendu}

Chaque fichier source doit inclure l'en-t\^ete de rendu d\'efini dans
\path{../../../resources/latex/deliverable-header.tex}, adapt\'e
\`a la syntaxe de commentaire du langage utilis\'e.

\section{Cibles Autoris\'ees Uniquement}

\begin{warning}
  Vous pouvez \textbf{uniquement} tester les cibles explicitement list\'ees
  dans l'exercice (labs, serveurs locaux, maquettes, binaires d\'esign\'es).
  Tester des syst\`emes de production r\'eels, des applications tierces
  ou toute cible qui ne vous appartient pas ou pour laquelle vous n'avez
  pas d'autorisation \'ecrite est interdit et peut entraîner une
  disqualification imm\'ediate.
\end{warning}

L'\'ethique du bug bounty s'applique tout au long de ce curriculum :
\begin{itemize}
  \item Aucun test sur des programmes en production sans un accord de
        p\'erim\`etre valide.
  \item Aucun d\'eni de service sauf si l'exercice l'exige explicitement.
  \item Aucune exfiltration de donn\'ees au-del\`a de ce qui est
        n\'ecessaire pour d\'emontrer la vuln\'erabilit\'e.
  \item Si vous acc\'edez accidentellement \`a des donn\'ees r\'eelles
        d'utilisateurs, arr\^etez imm\'ediatement et documentez ce qui
        s'est pass\'e.
\end{itemize}

\section{Citez Toutes Vos Sources}
\label{sec:cite-sources}

Toute ressource consult\'ee pendant votre travail doit \^etre cit\'ee.
Cela inclut :

\begin{itemize}
  \item Articles acad\'emiques et billets de blog.
  \item Rapports publics de bug bounty (HackerOne, Bugcrowd, etc.).
  \item Documentation officielle (RFC, pages de manuel, documentation
        fournisseur).
  \item \textbf{Outils d'IA} : nom du mod\`ele, r\'esum\'e de la
        question, date et m\'ethode de v\'erification (voir le chapitre
        Instructions IA pour le format de citation exact).
  \item Extraits de code adapt\'es de projets open-source (avec
        attribution de licence si requis).
\end{itemize}

\subsection*{Format de citation}

\begin{minted}{markdown}
# R\'ef\'erences

- [1] PortSwigger, "Business Logic Flaws",
  https://portswigger.net/web-security/logic-flaws
  (consult\'e le 2025-03-15).
- [2] Rapport HackerOne #123456, "IDOR on /api/users",
  anonymis\'e, consult\'e via divulgation publique.
- [3] G\'en\'er\'e avec l'assistance de Claude 3.5 Sonnet
  le 2025-03-14. Question : "Write a Python script that...".
  V\'erifi\'e par : revue manuelle de chaque requ\^ete HTTP.
\end{minted}

\subsection*{Attribution en ligne}

Si un exercice vous demande explicitement d'utiliser une technique ou
un outil sp\'ecifique, vous n'avez pas besoin de citer l'exercice
lui-m\^eme comme source. Citez tout ce qui se trouve en dehors de
l'\'enonc\'e de l'exercice.

\chapter{Instructions sur l'IA}

\section{Aperçu de la Politique}

Les outils de codage assistés par IA (ChatGPT, Claude, Copilot, Gemini,
Cursor, Codeium, et similaires) sont \textbf{autorisés} tout au long
de ce cursus. La recherche en sécurité est un domaine où la maîtrise
de l'IA devient essentielle --- apprendre à utiliser ces outils
efficacement fait partie de votre formation.

Cependant, l'objectif est la \textbf{compréhension}, pas l'automatisation.
Si vous ne pouvez pas expliquer chaque ligne de votre exploit sans
l'aide de l'IA, vous n'avez pas terminé l'exercice.

\begin{warning}
  Utiliser l'IA pour contourner l'objectif pédagogique d'un exercice
  --- par exemple, lui demander de rétro-ingéniérer un binaire pour
  ne pas avoir à apprendre le désassembleur, ou de générer une chaîne
  d'exploitation complète que vous ne comprenez pas --- va à l'encontre
  du but du cursus et sera traité comme une fraude académique.
\end{warning}

\section{Utilisations Autorisées}

\begin{itemize}
  \item \textbf{Explication de concepts}. Demander à l'IA de clarifier
        une classe de vulnérabilité, un sous-système du noyau, ou un
        protocole réseau.
  \item \textbf{Échafaudage de code}. Générer du code de départ pour
        un harnais, un fuzzer, ou un squelette de PoC --- puis le
        modifier et l'étendre vous-même.
  \item \textbf{Débogage}. Expliquer des messages d'erreur, la sortie
        d'un désassembleur, ou un comportement de crash inattendu.
  \item \textbf{Résumé de documentation}. Condenser une longue RFC ou
        un article technique en notes exploitables.
  \item \textbf{Traduction}. Traduire du contenu technique entre
        l'anglais et le français (les deux langues sont utilisées
        dans ce cursus).
  \item \textbf{Rédaction de rapports}. Demander de l'aide pour
        structurer un rapport de bug bounty --- mais remplir le
        contenu technique vous-même.
\end{itemize}

\section{Utilisations Interdites}

\begin{itemize}
  \item \textbf{Génération complète d'exploit}. Demander à l'IA
        d'écrire l'exploit complet d'un exercice, puis le soumettre
        sans le comprendre.
  \item \textbf{Contournement des contraintes de l'exercice}. Utiliser
        l'IA pour contourner une restriction que l'exercice impose
        explicitement (p.\ ex.\ «~n'utilisez pas de scanners
        automatisés~»).
  \item \textbf{Découverte automatisée de vulnérabilités}. Exécuter des
        outils de scan pilotés par IA qui remplacent l'analyse manuelle
        que vous êtes censé pratiquer.
  \item \textbf{Rédaction des grilles d'évaluation}. Vous n'êtes pas
        autorisé à utiliser l'IA pour générer des critères d'évaluation
        par les pairs ou des réponses.
  \item \textbf{Plagiat sans attribution}. Copier le résultat de l'IA
        sans le citer (voir Section~\ref{sec:ai-citation}).
\end{itemize}

\section{Responsabilité Personnelle}

Lorsque vous demandez à l'IA de générer du code d'exploit, vous
devenez responsable de chaque octet de ce résultat avant de
l'utiliser contre une cible.

\begin{enumerate}
  \item \textbf{Lisez le code généré} ligne par ligne.
  \item \textbf{Identifiez toutes les hypothèses} faites par l'IA
        (architecture, version du noyau, drapeaux du compilateur,
        topologie réseau).
  \item \textbf{Testez en isolation} avant d'exécuter contre une
        cible d'exercice.
  \item \textbf{Comprenez les modes d'échec}. Si l'exploit généré
        plante ou se bloque, vous devez diagnostiquer et corriger
        le problème vous-même.
\end{enumerate}

\hint{Si l'IA génère du code que vous ne comprenez pas, demandez-lui
de l'expliquer. Si vous ne comprenez toujours pas après
l'explication, le code n'est pas prêt à être utilisé. Passez à
une approche plus simple que vous maîtrisez entièrement.}

\section{Citation Obligatoire de l'IA}
\label{sec:ai-citation}

Chaque fichier contenant du code généré ou assisté par l'IA doit
inclure un commentaire de citation. Le format est~:

\begin{minted}{python}
# Generated with assistance from <tool> on <date>.
# Prompt: <one-line summary of the prompt>
# Verified by: <method used to verify the output>
\end{minted}

\subsection*{Exemples}

\textbf{Python}~:
\begin{minted}{python}
# Generated with assistance from Claude 3.5 Sonnet on 2025-03-14.
# Prompt: "Write a Python script that sends HTTP requests with
# modified price fields to /api/checkout/pay"
# Verified by: manual review of each request and response
\end{minted}

\textbf{C}~:
\begin{minted}{c}
/* Generated with assistance from GitHub Copilot on 2025-03-14.
 * Prompt: "Implement a stack-based buffer overflow with
 * ret2libc for x86_64 Linux"
 * Verified by: compiled with -fno-stack-protector and
 * tested against local binary with ASLR disabled */
\end{minted}

\textbf{Markdown (rapport)}~:
\begin{minted}{html}
<!--
Generated with assistance from ChatGPT-4 on 2025-03-14.
Prompt: "Structure a bug bounty report for IDOR with
sections: summary, steps, impact, mitigation"
Verified by: all technical content written manually,
AI only suggested structure
-->
\end{minted}

\subsection*{Interactions multiples avec l'IA}

Si vous avez utilisé plusieurs outils d'IA, ou itéré à travers
plusieurs invites, citez chaque interaction significative. Un
groupe d'invites connexes peut être résumé~:

\begin{minted}{html}
<!--
AI assistance used in this file:
- Claude 3.5 Sonnet (2025-03-14): harness skeleton generation
- ChatGPT-4 (2025-03-14): crash triage help
- Copilot (2025-03-15): inline autocomplete for loop logic
All generated code reviewed and modified manually.
-->
\end{minted}

\section{Évaluation par les Pairs et l'IA}

Lors de l'évaluation par les pairs, l'évaluateur peut~:

\begin{itemize}
  \item Vous demander d'\textbf{expliquer tout code généré par l'IA}
        ligne par ligne.
  \item Vous demander de \textbf{reproduire le raisonnement} sans
        l'aide de l'IA.
  \item Demander à voir votre \textbf{historique d'invites} pour un
        exercice donné.
  \item Déduire des points si la citation de l'IA est manquante ou
        si vous ne pouvez pas expliquer le code que l'IA a généré.
\end{itemize}

\begin{warning}
  Présenter le code généré par l'IA comme votre propre travail
  (défaut de citation) est du plagiat. Si l'évaluateur soupçonne
  du plagiat, il peut saisir l'équipe pédagogique pour examen.
\end{warning}

\section{Rappel Éthique}

\begin{itemize}
  \item Les modèles d'IA peuvent produire du code incorrect,
        dangereux, ou illégal. Vous êtes \textbf{responsable} de
        vérifier chaque résultat avant exécution.
  \item Les modèles d'IA entraînés sur la recherche publique en
        sécurité peuvent reproduire textuellement du code provenant
        de projets sous licence GPL. Vérifiez les licences avant
        d'intégrer ces éléments dans votre travail.
  \item Ne saisissez pas de code sensible ou propriétaire dans des
        outils d'IA hébergés dans le cloud. Utilisez des modèles
        locaux (Llama, Mistral, DeepSeek) pour les travaux impliquant
        des cibles non publiques.
  \item L'éthique de la recherche en sécurité que vous apprenez ici
        s'applique à l'utilisation de l'IA~: ne demandez pas à l'IA
        de générer des malwares, des rançongiciels, ou du code destiné
        à nuire à des systèmes en dehors du cadre de l'exercice.
\end{itemize}

\info{La maîtrise des outils d'IA est une compétence recherchée dans
la chasse aux bogues moderne. Utilisez-les sagement, citez-les
honnêtement, et conservez toujours votre propre compréhension de
chaque exploit que vous déployez.}


\chapter{Avant-propos}

Le papier ne dit pas « ce CVE est un bug de logique ». Il dit qu'un
composant a une \emph{forme architecturale}, et que cette forme prédit la
manière dont les correctifs vont échouer.

Trois familles suffisent pour commencer :

\begin{center}
\begin{tabularx}{\linewidth}{|p{0.22\linewidth}|X|X|}
\hline
\textbf{Famille} & \textbf{Question diagnostique} & \textbf{Mode d'échec typique} \\ \hline
Logical Engine & Le composant gère-t-il des objets avec un cycle de vie
explicite (create, destroy, abort, rollback) ? & Transition d'état
asymétrique : le chemin d'abandon ne préserve pas les invariants du chemin
nominal. \\ \hline
Predictive Oracle & Le composant analyse-t-il une entrée avant de
l'exécuter (vérificateur, JIT, type checker) ? & L'analyse est incomplète :
une passe croit une propriété que le backend ne préserve pas. \\ \hline
Concurrent Proxy & Le composant partage-t-il de la mémoire ou des
ressources avec une synchronisation minimale (rings, files, completions) ? &
Le moment de la validation n'est pas le moment de l'utilisation. \\ \hline
\end{tabularx}
\end{center}

\begin{info}
Une famille décrit une \emph{forme de gestion}, pas la gravité du bug. Un
composant peut cumuler plusieurs familles ; choisissez celle dont le mode
d'échec explique le mieux \emph{pourquoi le patch a échoué}.
\end{info}

\chapter{Aperçu du module}

\begin{center}
\begin{tabularx}{\linewidth}{|c|l|c|c|}
\hline
\textbf{Ex} & \textbf{Titre} & \textbf{Difficulté} & \textbf{Points} \\ \hline
00 & Classer six composants & Facile & 3 \\
01 & Chemins nominaux et chemins d'abandon & Moyen & 4 \\
02 & Cartographier huit CVE & Moyen & 5 \\
03 & Écrire une fiche de famille & Difficile & 4 \\ \hline
& \textbf{Total} & & \textbf{16} \\ \hline
\end{tabularx}
\end{center}

\startexercices

\chapter{Exercise 00}
\extitle{Classer six composants}
\turnindir{ex00}
\exnumber{00}
\exfiles{family-classification.md}
\exauthorize{un éditeur de texte, accès aux sources}
\makeheaderfiles

Classez les six composants suivants dans une famille principale. Pour
chacun, justifiez en trois phrases maximum et indiquez la question qui a
tranché.

\begin{enumerate}
  \item nf\_tables (\texttt{net/netfilter/nf\_tables\_api.c})
  \item Le vérificateur eBPF (\texttt{kernel/bpf/verifier.c})
  \item io\_uring (\texttt{io\_uring/})
  \item La chaîne d'authentification Windows (NTLM, Kerberos, SMB)
  \item La normalisation de chemins dans Node.js
  \item La désérialisation dans Apache MINA
\end{enumerate}

Remplissez \texttt{family-classification.md} :

\begin{center}
\begin{tabularx}{\linewidth}{|p{0.18\linewidth}|p{0.18\linewidth}|X|p{0.2\linewidth}|}
\hline
\textbf{Composant} & \textbf{Famille} & \textbf{Justification} & \textbf{Question décisive} \\ \hline
 & & & \\ \hline
 & & & \\ \hline
 & & & \\ \hline
 & & & \\ \hline
 & & & \\ \hline
 & & & \\ \hline
\end{tabularx}
\end{center}

\hint{Ne classez pas en fonction du langage ou du sous-système. Un
vérificateur écrit en C reste un Predictive Oracle. Une machine à états
d'authentification reste un Logical Engine.}

\chapter{Exercise 01}
\extitle{Chemins nominaux et chemins d'abandon}
\turnindir{ex01}
\exnumber{01}
\exfiles{state-machine.md}
\exauthorize{un éditeur de texte, Python 3}
\makeheaderfiles

Voici une machine à états simplifiée. Elle crée un objet, peut l'abandonner
en cours de route, puis le détruit.

\begin{console}
state = "NEW"
refcount = 1

def create():
    global state, refcount
    state = "ACTIVE"

def abort():
    global state
    state = "ABORTED"
    # le refcount n'est pas decremente

def destroy():
    global state, refcount
    if state == "ACTIVE":
        refcount -= 1
    state = "DESTROYED"
\end{console}

\subsection*{Mission}

Dans \texttt{state-machine.md} :

\begin{enumerate}
  \item Décrivez le chemin nominal : \texttt{NEW -> ACTIVE -> DESTROYED}.
  \item Décrivez le chemin d'abandon : \texttt{NEW -> ABORTED -> DESTROYED}.
  \item Quel invariant est violé par \texttt{abort()} ?
  \item Écrivez un scénario de trois appels qui produit un refcount
        incohérent.
  \item Proposez un correctif minimal et expliquez quel test manquait.
\end{enumerate}

\begin{warn}
Ce code est un jouet pédagogique. Ne le transposez pas directement à un
allocateur noyau : les objets réels ont des règles de verrouillage, de
génération et de RCU qui rendent l'analyse plus subtile.
\end{warn}

\chapter{Exercise 02}
\extitle{Cartographier huit CVE}
\turnindir{ex02}
\exnumber{02}
\exfiles{cve-family-map.md}
\exauthorize{un éditeur de texte, NVD, papier Phrack}
\makeheaderfiles

Remplissez \texttt{cve-family-map.md} pour les CVE suivants. Pour chaque
ligne : famille principale, hypothèse du correctif, mode de mutation
(\emph{incomplete-fix variant}, \emph{architectural scatter}, ou
\emph{correctif complet}).

\begin{center}
\begin{tabularx}{\linewidth}{|p{0.2\linewidth}|X|X|}
\hline
\textbf{CVE} & \textbf{Hypothèse du correctif} & \textbf{Mode de mutation} \\ \hline
CVE-2023-35001 & & \\ \hline
CVE-2024-1086 & & \\ \hline
CVE-2023-4244 & & \\ \hline
CVE-2024-1085 & & \\ \hline
CVE-2020-27194 & & \\ \hline
CVE-2021-3490 & & \\ \hline
CVE-2023-2598 & & \\ \hline
CVE-2024-0582 & & \\ \hline
\end{tabularx}
\end{center}

Ajoutez deux paragraphes :

\begin{itemize}
  \item Pourquoi CVE-2023-35001 et CVE-2024-1086 sont-ils une paire
        d'\emph{architectural scatter} et non un enchaînement de correctifs
        incomplets ?
  \item Pourquoi CVE-2023-4244 et CVE-2024-1085 sont-ils, eux, un cas de
        correctif incomplet ?
\end{itemize}

\hint{La distinction ne dépend pas de la proximité des dates. Elle dépend de
la question : le second bug était-il déjà là avant le premier correctif, ou
a-t-il été rendu atteignable par la suppression d'une garde ?}

\chapter{Exercise 03}
\extitle{Écrire une fiche de famille}
\turnindir{ex03}
\exnumber{03}
\exfiles{family-card.md}
\exauthorize{un éditeur de texte}
\makeheaderfiles

Choisissez une famille et un composant. Produisez une fiche opérationnelle
de deux pages maximum, destinée à un auditeur qui découvre le composant.

La fiche doit contenir :

\begin{enumerate}
  \item \textbf{Forme} : les objets, états ou ressources gérés.
  \item \textbf{Chemin nominal} : la séquence heureuse.
  \item \textbf{Chemins d'exception} : abort, rollback, cancel, timeout,
        concurrence.
  \item \textbf{Invariants} : ce qui doit rester vrai dans tous les chemins.
  \item \textbf{Seams} : les fonctions ou champs à regarder en premier.
  \item \textbf{Cinq questions} : des questions de chasse, pas des
        généralités.
  \item \textbf{Anti-pattern} : une erreur d'analyse fréquente sur cette
        famille.
\end{enumerate}

\chapter{Rendu et évaluation}

\begin{console}
M-EPI-001-mutation-model/
|-- ex00/family-classification.md
|-- ex01/state-machine.md
|-- ex02/cve-family-map.md
+-- ex03/family-card.md
\end{console}

\begin{center}
\begin{tabularx}{\linewidth}{|l|X|c|}
\hline
\textbf{Critère} & \textbf{Attendu} & \textbf{Points} \\ \hline
Classification & Famille correcte et question décisive explicite & 4 \\ \hline
Chemins & Nominal et abandon distingués, invariant nommé & 4 \\ \hline
CVE & Mode de mutation justifié, pas seulement étiqueté & 4 \\ \hline
Fiche & Utilisable par un autre auditeur sans contexte oral & 4 \\ \hline
\end{tabularx}
\end{center}

\chapter{Ressources par exercice}

\subsection*{Exercice 00 --- Classer six composants}
\begin{itemize}
  \item Papier, section 3 : les trois questions diagnostiques.
  \item Dépôt compagnon : \texttt{validation/methodology.md}.
  \item NVD : \url{https://nvd.nist.gov/}
  \item À chercher : \texttt{Logical Engine}, \texttt{Predictive Oracle},
        \texttt{Concurrent Proxy}, puis le cycle de vie de chaque composant.
\end{itemize}

\subsection*{Exercice 01 --- Chemins nominaux et d'abandon}
\begin{itemize}
  \item Papier, section 3 : transitions d'état asymétriques.
  \item Documentation RCU : \url{https://docs.kernel.org/RCU/}
  \item À chercher : \texttt{abort}, \texttt{rollback}, \texttt{cancel},
        \texttt{refcount}, \texttt{state machine}.
\end{itemize}

\subsection*{Exercice 02 --- Cartographier huit CVE}
\begin{itemize}
  \item NVD : CVE-2023-35001, CVE-2024-1086, CVE-2023-4244,
        CVE-2024-1085, CVE-2020-27194, CVE-2021-3490, CVE-2023-2598,
        CVE-2024-0582.
  \item Papier, sections 6.2, 6.3, 14 et 15.
  \item Dépôt compagnon : \texttt{validation/dataset.csv}.
  \item À chercher : \texttt{architectural scatter},
        \texttt{incomplete-fix variant}, \texttt{structural symmetry}.
\end{itemize}

\subsection*{Exercice 03 --- Écrire une fiche de famille}
\begin{itemize}
  \item Papier, sections 2 à 4 et 18.7.
  \item Dépôt compagnon : \texttt{rules/TRIAGE.md}.
  \item NVD : \url{https://nvd.nist.gov/}
  \item À chercher : \texttt{seam}, \texttt{invariant}, \texttt{nominal path},
        \texttt{exception path}.
\end{itemize}

\chapter{Sources}
\begin{itemize}
  \item Phrack, \emph{Security is failure: Every Patch Is a Guess},
        sections 2, 3, 6.2, 6.3, 14 et 15.
  \item Dépôt compagnon :
        \texttt{validation/methodology.md} et \texttt{validation/dataset.csv}.
  \item NVD : \url{https://nvd.nist.gov/}
\end{itemize}

\end{document}
